\documentclass[11pt]{article}
\usepackage[margin=1in]{geometry}
\usepackage[T1]{fontenc}
\usepackage[utf8]{inputenc}
\usepackage{lmodern}
\usepackage{hyperref}
\usepackage{enumitem}
\title{Concept Notes: A Tensegrity-Inspired Humanoid with Textile Outer Shell}
\author{Working notes}
\date{\today}

\begin{document}
\maketitle

\section*{Thesis}
A humanoid can be built as a hybrid of:
\begin{enumerate}[leftmargin=*]
  \item a tensegrity-inspired internal load-bearing network,
  \item discontinuous stiff inserts (``bones'') captured by garment pockets,
  \item tendon transmission routed through textile channels, and
  \item electric actuation (BLDC + gearbox + spool, and/or linear actuators) located in low-risk body zones.
\end{enumerate}

This architecture is plausible and technically interesting because it can combine impact tolerance, human-safe compliance, modular repair, and apparel-style manufacturing scalability.

\section*{Why this is compelling}
\begin{itemize}[leftmargin=*]
  \item \textbf{Safety and compliance}: tensegrity and tendon elasticity reduce peak contact forces.
  \item \textbf{Weight distribution}: actuator packs can be concentrated near torso/pelvis, while distal limbs remain lightweight.
  \item \textbf{Manufacturing path}: cut-sew-laminate, reinforcement tapes, and pocketed inserts are familiar to apparel and soft goods production.
  \item \textbf{Serviceability}: shell panels and tendon modules can be replaced without disassembling a monolithic rigid frame.
\end{itemize}

\section*{Core architecture proposal}
\subsection*{1) Structural layers}
\begin{enumerate}[leftmargin=*]
  \item \textbf{Inner skeleton}: sparse compression members and tension network (tensegrity-inspired, not necessarily pure tensegrity at whole-body scale).
  \item \textbf{Bone inserts}: carbon, glass fiber, or printed stiffeners in indexed garment pockets.
  \item \textbf{Textile shell}: multilayer technical fabric with:
  \begin{itemize}[leftmargin=*]
    \item tendon routing sleeves,
    \item local abrasion liners,
    \item anisotropic reinforcement along principal load paths.
  \end{itemize}
\end{enumerate}

\subsection*{2) Actuation options}
\begin{itemize}[leftmargin=*]
  \item \textbf{BLDC + reduction + spool}:
  high efficiency and controllability, but requires careful spool geometry and cable management.
  \item \textbf{Linear actuators}:
  simple kinematics at selected joints, typically heavier for the same peak force.
  \item \textbf{Hybrid}:
  rotary tendon actuation for major joints, short-stroke linear units for stiffness tuning or pretension.
\end{itemize}

\subsection*{3) Tendon transmission details}
\begin{itemize}[leftmargin=*]
  \item Use antagonistic tendon pairs for bidirectional control and stiffness modulation.
  \item Include inline compliance and pretension adjusters.
  \item Model friction and hysteresis explicitly (Bowden effects, bend-radius penalties, creep).
\end{itemize}

\section*{Main technical risks}
\begin{enumerate}[leftmargin=*]
  \item \textbf{Tendon friction/hysteresis}: can dominate precision and bandwidth limits.
  \item \textbf{Textile fatigue}: seam creep, laminate delamination, and abrasion at routing bends.
  \item \textbf{Calibration drift}: tension distribution changes over time and with temperature/humidity.
  \item \textbf{Control complexity}: high-DOF compliant systems need robust estimation and model adaptation.
\end{enumerate}

\section*{Practical design rules (first prototype)}
\begin{enumerate}[leftmargin=*]
  \item Start with \textbf{one limb module} (for example, shoulder-elbow) before full humanoid integration.
  \item Keep \textbf{actuator modules external and swappable} in revision A.
  \item Build \textbf{instrumented tendon paths}: tension, displacement, and bend-state sensing.
  \item Add \textbf{hard end-stops and passive fail-safe postures} for every powered DOF.
  \item Define a \textbf{fatigue test plan early}: repeated cycles at realistic duty and environmental conditions.
\end{enumerate}

\section*{Suggested phased roadmap}
\subsection*{Phase 0: Bench characterization}
Measure tendon friction loops, spool repeatability, textile channel wear, and stiffness range of one joint module.

\subsection*{Phase 1: Upper-body demonstrator}
Build a torso + one arm with pocketed inserts and textile shell; validate force output, compliance envelopes, and impact behavior.

\subsection*{Phase 2: Mobility integration}
Add pelvis and lower-limb modules; evaluate quasi-static gait and fall-tolerant recovery behaviors.

\subsection*{Phase 3: Reliability and manufacturing}
Transition from hand-built prototypes to reproducible cut-sew-laminate process with QC metrics on preload and tendon routing.

\section*{Positioning against prior work}
Your concept is not just ``another humanoid'' and not just ``another soft robot.''
It sits at a valuable intersection:
\begin{itemize}[leftmargin=*]
  \item tensegrity-inspired resilience,
  \item anthropomimetic tendon-driven actuation,
  \item and textile-first manufacturability.
\end{itemize}

That intersection is underexplored and could produce both strong research contributions and practical product pathways.

\section*{NSF FRR-oriented preliminary design}
\subsection*{Working objective}
Demonstrate a \textbf{textile-structured, tendon-driven, compliant arm module} that can:
\begin{enumerate}[leftmargin=*]
  \item regulate endpoint force under disturbances,
  \item vary apparent joint stiffness through antagonistic co-contraction, and
  \item survive repeated cyclic loading with bounded performance drift.
\end{enumerate}

\subsection*{Preliminary module definition (Design A0)}
\begin{itemize}[leftmargin=*]
  \item \textbf{Kinematics}: 2 rotational DOF (shoulder pitch, elbow flexion).
  \item \textbf{Structure}: one rigid torso anchor, two compliant limb segments, indexed pockets for stiff inserts.
  \item \textbf{Transmission}: 2 antagonistic tendon pairs (one pair per DOF).
  \item \textbf{Actuation}: BLDC + gearbox + spool modules mounted near the torso anchor.
  \item \textbf{Sensing}: joint position, tendon displacement, tendon tension, motor current.
  \item \textbf{Safety}: passive end stops, current limits, and low-stiffness startup mode.
\end{itemize}

\subsection*{FRR-style testable hypotheses}
\begin{itemize}[leftmargin=*]
  \item \textbf{H1}: antagonistic co-contraction increases closed-loop disturbance rejection without requiring rigid links.
  \item \textbf{H2}: textile-integrated routing can maintain repeatable kinematics over useful cycle counts when reinforcement and bend-radius constraints are enforced.
  \item \textbf{H3}: a reduced-order simulation model can predict stiffness and range-of-motion trends well enough to drive design iterations.
\end{itemize}

\section*{MuJoCo modeling plan (immediate next step)}
\subsection*{Modeling scope}
Start with a reduced model that captures the most important design variables:
\begin{itemize}[leftmargin=*]
  \item 2-DOF arm kinematics,
  \item tendon-driven actuation mapped to each joint,
  \item damping/inertia/friction abstractions,
  \item co-contraction commands for variable apparent stiffness.
\end{itemize}

\subsection*{Minimum simulation outputs}
\begin{enumerate}[leftmargin=*]
  \item Range of motion envelopes.
  \item Step response and settling time at both joints.
  \item Endpoint displacement under impulse or constant external force.
  \item Effective stiffness versus co-contraction command.
\end{enumerate}

\subsection*{Design iteration loop}
\begin{enumerate}[leftmargin=*]
  \item Tune masses, damping, and actuator gains for realistic baseline behavior.
  \item Sweep tendon and control parameters.
  \item Select top candidate configurations by stability and force-tracking metrics.
  \item Export candidate parameter set for hardware bench prototype.
\end{enumerate}

\subsection*{Near-term milestone}
Within one cycle, deliver:
\begin{itemize}[leftmargin=*]
  \item a baseline MuJoCo model of the shoulder-elbow module,
  \item a scripted experiment set (ROM, disturbance, co-contraction sweep), and
  \item a parameter table that directly informs the first physical prototype build.
\end{itemize}

\section*{References in this workspace}
See:
\begin{itemize}[leftmargin=*]
  \item \texttt{../references/related\_work.md}
  \item \texttt{../references/library.bib}
\end{itemize}

\end{document}
